%==============================================================================
%  Apéndice 1: Implicaciones ambientales y de desarrollo sostenible.
%  Contenido compartido por el informe (Plantilla_Informe_TFG.tex) y por su copia
%  aparte (Apendice1_Implicaciones.tex). El título lo pone cada documento.
%==============================================================================

\textit{AutoPipe} es una librería de software de uso abierto que automatiza una decisión de
ingeniería habitualmente manual: qué preprocesamiento aplicar a los datos antes de entrenar
un modelo. Este apéndice analiza sus implicaciones en seis aspectos, y relaciona cada uno con
los Objetivos de Desarrollo Sostenible (ODS) de la Agenda 2030 de las Naciones Unidas
\cite{un2015agenda}. El análisis distingue lo que el proyecto hizo y midió de lo que su uso
podría implicar: una librería de propósito general hereda las consecuencias de los datos y
de los modelos sobre los que se aplique, y parte de esas consecuencias queda en manos de
quien la usa.

\subsubsection*{Aspecto social}

\textbf{Identificación y análisis.} Elegir el preprocesamiento exige experiencia
especializada y tiempo de prueba y error. Automatizarlo, con documentación en español y
licencia abierta, reduce la barrera de entrada para equipos pequeños, estudiantes e
investigadores sin un especialista en aprendizaje automático. La automatización tiene un
riesgo social propio: ocultar decisiones que afectan al resultado. El proyecto lo mitiga
exponiendo la tabla completa de candidatos y un reporte reproducible de cada búsqueda, y sus
propios resultados muestran por qué no basta con confiar en una recomendación: la estrategia
por similitud, la más intuitiva, fue la menos eficaz de las evaluadas.

\textbf{Relación con los ODS.} Se relaciona con el ODS~4, educación de calidad, por el
acceso abierto a una herramienta documentada y a un protocolo experimental reutilizable;
con el ODS~9, industria, innovación e infraestructura, por la difusión de capacidades de
aprendizaje automático; y con el ODS~10, reducción de las desigualdades, por disminuir la
dependencia de recursos especializados.

\textbf{Proyecciones futuras.} La automatización de etapas del aprendizaje automático
ampliará quién puede usarlo. El ingeniero en computadores tendrá una responsabilidad
creciente en que esas herramientas expliquen sus decisiones y en formar a quienes las usan
para interpretar sus resultados en lugar de aceptarlos.

\subsubsection*{Aspecto económico}

\textbf{Identificación y análisis.} El costo de seleccionar un preprocesamiento es tiempo de
cómputo y tiempo de especialistas. El arranque en caliente redujo en el segundo experimento
un 35,8\,\% la mediana del \textit{regret} normalizado, que mide cuánto se aleja la búsqueda
del óptimo a lo largo de su presupuesto: para una misma calidad hacen falta menos
evaluaciones. El proyecto se desarrolló íntegramente con herramientas libres y datos
públicos, sin costos de licencias, y su producto se distribuye bajo licencia MIT.

\textbf{Relación con los ODS.} Contribuye al ODS~8, trabajo decente y crecimiento económico,
por el aumento de productividad en tareas técnicas, y al ODS~9 por facilitar la innovación
en organizaciones con recursos de cómputo limitados.

\textbf{Proyecciones futuras.} Las organizaciones pequeñas y las universidades con
presupuestos acotados son las que más se benefician de reducir evaluaciones. El valor del
trabajo técnico se desplazará desde ejecutar búsquedas hacia validar y auditar sus
resultados.

\subsubsection*{Aspecto ambiental}

\textbf{Identificación y análisis.} El cómputo tiene una huella energética, y la literatura
ha mostrado que la del aprendizaje automático puede ser considerable y que conviene
reportarla junto a los resultados \cite{strubell2019energy, schwartz2020green}. Los
experimentos de este proyecto se ejecutaron en hardware local, sin aceleradores gráficos ni
infraestructura en la nube: el primer experimento tomó 13,4 minutos, el cálculo de los
óptimos de referencia unos 62 minutos, el segundo experimento 133 minutos y el análisis de
sensibilidad 90 minutos. El propósito mismo del arranque en caliente tiene una lectura
ambiental, porque reutiliza lo aprendido en búsquedas anteriores en lugar de repetirlas.

\textbf{Relación con los ODS.} Se relaciona con el ODS~12, producción y consumo responsables,
por el uso eficiente de recursos de cómputo, y con el ODS~13, acción por el clima, por la
contribución indirecta a reducir el consumo energético asociado a la experimentación.

\textbf{Proyecciones futuras.} Compartir conocimiento entre búsquedas, a escala de muchos
usuarios, podría evitar cómputo redundante. Para que ese beneficio sea verificable, las
herramientas de AutoML deberían reportar el costo computacional de cada búsqueda, como ya
propone la literatura.

\subsubsection*{Aspecto legal}

\textbf{Identificación y análisis.} El proyecto combina tres regímenes de licencia. El código
se publica bajo licencia MIT. Las imágenes de MVTec AD \cite{bergmann2019mvtec} se usan bajo
la licencia CC BY-NC-SA 4.0, que prohíbe el uso comercial y condiciona su redistribución; el
proyecto optó por no redistribuirlas: no se versionan en el repositorio, y una prueba falla si
los reportes que contienen miniaturas dejaran de excluirse. Los conjuntos tabulares provienen de la suite OpenML-CC18
\cite{openml-cc18}, cuya licencia se verificó para la suite y no conjunto por conjunto,
limitación declarada en los criterios de inclusión del corpus. La librería no recolecta datos
personales, pero puede aplicarse sobre ellos: en Costa Rica, la responsabilidad sobre ese
tratamiento corresponde a quien lo realice, conforme a la Ley N.\textsuperscript{o}~8968
\cite{ley8968}.

\textbf{Relación con los ODS.} Se vincula con el ODS~16, paz, justicia e instituciones sólidas,
en lo relativo a la transparencia y la rendición de cuentas: los pre-registros congelados y
los reportes reproducibles permiten auditar cómo se obtuvo cada resultado.

\textbf{Proyecciones futuras.} Las regulaciones sobre inteligencia artificial tienden a exigir
trazabilidad y documentación de los sistemas automatizados. Prácticas como las de este
proyecto ---versionar el protocolo antes de ejecutar y regenerar los resultados desde sus
datos--- facilitarán ese cumplimiento, del mismo modo que las licencias de los conjuntos de
datos de referencia exigirán una verificación individual cada vez más rigurosa.

\subsubsection*{Aspecto de seguridad}

\textbf{Identificación y análisis.} Tres riesgos de seguridad y confiabilidad se abordaron
por diseño. La fuga de información entre entrenamiento y validación, que produce resultados
optimistas y engañosos, se impide en el protocolo y se vigila con pruebas escritas para
fallar. Las configuraciones que se inyectan desde una base de meta-conocimiento se validan
antes de ejecutarse, y las bases y configuraciones se leen con formatos de datos ---JSON y
YAML en modo seguro--- que no ejecutan código. Y los fallos silenciosos se hicieron
explícitos: cuando ningún candidato se puede evaluar, la librería lo avisa. Queda un riesgo
de uso que conviene señalar: en la detección de anomalías en imágenes el modelo base obtuvo
resultados cercanos al azar en una de las categorías, de modo que los resultados de este
proyecto no son aptos para inspección industrial crítica sin un modelo más capaz.

\textbf{Relación con los ODS.} Se relaciona con el ODS~9, por la construcción de
infraestructura tecnológica fiable, y con el ODS~8 en lo relativo a entornos de trabajo
seguros, donde la inspección automatizada tiene aplicación.

\textbf{Proyecciones futuras.} A medida que se compartan bases de conocimiento y
configuraciones entre organizaciones, su integridad será un problema de seguridad de la
cadena de suministro del software, que exigirá versionado, validación y firma de esos
artefactos.

\subsubsection*{Aspecto de salud}

\textbf{Identificación y análisis.} El proyecto no trabajó con pacientes ni con datos
clínicos propios, pero su corpus público incluye conjuntos de origen médico, como los de
diabetes, pacientes hepáticos, cáncer de mama y afecciones tiroideas. Eso muestra que la
librería puede aplicarse a datos de salud, y obliga a precisar el alcance de sus resultados:
seleccionar el preprocesamiento que maximiza una métrica sobre una partición de validación es
un resultado de referencia, no una validación clínica. El uso de métricas macro, que dan el
mismo peso a todas las clases, es pertinente en ese contexto, donde la clase de interés suele
ser la minoritaria.

\textbf{Relación con los ODS.} Se relaciona con el ODS~3, salud y bienestar, por el potencial
de las técnicas de aprendizaje automático en el apoyo al diagnóstico y la detección temprana.

\textbf{Proyecciones futuras.} La aplicación de herramientas de AutoML a datos clínicos
requerirá validación externa, análisis de sesgos entre poblaciones y aprobación regulatoria.
La cultura de pre-registro y reproducibilidad que este proyecto aplicó es la misma que exige
la investigación clínica, y es un puente natural entre ambas prácticas.

